% Section: Abstract
% Source: paper_draft_v0.1.md § Abstract

Cyber-hygiene anomaly detection---identifying stale privileged accounts, dormant
account reactivations, endpoint coverage gaps, patch noncompliance clusters, and
telemetry missingness---is operationally important but poorly served by existing
public benchmarks, which focus predominantly on network intrusion or
attack-emulation telemetry. We introduce \hygienebench, an open, reproducible
benchmark that jointly covers Active Directory identity state, endpoint patch
posture, vulnerability exposure, and telemetry freshness across seven evaluation
tasks (T1--T7) and twelve anomaly classes. \hygienebench is built entirely from
synthetic, seeded telemetry generated from publicly citable structural priors
(NIST NVD severity distributions, Verizon DBIR 2026 patch-lag
aggregates~\cite{verizon_dbir_2026}, CISA BOD 23-01 telemetry-cadence
requirements), with no employer or production data at any stage.

We evaluate eight methods---a task-specific rule baseline (M1), a hybrid risk
scorer (M2), Isolation Forest (M3), Local Outlier Factor (M4), One-Class SVM
(M5), a linear autoencoder (M6), population-level temporal z-score (M7), and a
graph-augmented Isolation Forest (M8)---across five telemetry conditions
(baseline, fresh, stale, missing, unsupervised). Applying a pre-registered failure
protocol ($\Delta<0.05$ AP and $\Delta<0.05$ P@k in $\geq$2/3 seeds), we find that
\textbf{83.8\% of (condition, task, method) configurations fail to improve on the
rule baseline}, with ML adding meaningful signal on T2 (group-membership-drift
detection, best ML: M7 temporal z-score, $+$0.185 AP over rule) and T5
(patch-vulnerability hygiene, best ML: M5 OCSVM, $+$0.210 AP over rule).
Telemetry staleness (C-STALE) consistently degrades detection performance relative
to the baseline condition. We release the generator, datasets, and evaluation
harness to enable reproducible comparison of future methods.
